% from ega4-2-fr.tex
\oldpage[IV-2]{158}
Soit $Y'=Y_{\mathrm{red}}$, qui est un sous-préschéma fermé intègre de $Y$
ayant même espace sous-jacent, et défini par un Idéal localement nilpotent
$\mathcal J$ de $\mathcal O_Y$ (I, 6.1.6). Il en résulte que
$\operatorname{gr}_{\mathcal J}^{\bullet}(\mathcal O_Y)$ est un
$\mathcal O_{Y'}$-Module cohérent, et comme $Y'$ est intègre, il y a un
voisinage ouvert $V$ de $y$ dans $Y'$ tel que
$\operatorname{gr}_{\mathcal J}^{\bullet}(\mathcal O_Y)|V$ soit localement
libre (0$_{\rm I}$, 5.2.7); autrement dit, quitte à remplacer $X$ par un
voisinage de $y$, on peut supposer que $\mathcal O_Y$ est normalement plat
le long de $Y'$; on déduit donc de (6.10.3) que l'on a
\[
\dim(\mathcal O_{Y,x})=\dim(\mathcal O_{Y',x}),\qquad
\operatorname{prof}(\mathcal O_{Y,x})=
\operatorname{prof}(\mathcal O_{Y',x})
\]
pour tout $x\in Y$; cela permet de se borner au cas où le sous-préschéma
fermé $Y$ est intègre.

Cela étant, en vertu de (6.10.2), on peut, en remplaçant $X$ par un
voisinage ouvert de $y$, supposer que $\mathcal F$ est normalement plat le
long de $Y$. Comme $y\in\operatorname{Supp}(\mathcal F)$, la relation
(6.10.6.1) résulte de (6.10.5.2) et (6.10.5.4). Posons maintenant
$p=\operatorname{prof}(\mathcal F_y)$; remplaçant au besoin $X$ par un
voisinage ouvert de $y$, on peut supposer qu'il existe $p$ sections $f_i$
$(1\leq i\leq p)$ de $\mathcal O_X$ au-dessus de $X$, formant une suite
$\mathcal F$-régulière, tels que les $(f_i)_y$ appartiennent à l'idéal
maximal $\mathfrak m_y$ de $\mathcal O_{X,y}$ et que
\[
\operatorname{prof}\left(\mathcal F_y/
\left(\sum_{i=1}^{p}(f_i)_y\mathcal F_y\right)\right)=0
\]
(0, 16.4.6). Si $\mathcal J$ est l'Idéal de $\mathcal O_X$ qui définit
$Y$, on a $\mathcal J_y=\mathfrak m_y$ et, en remplaçant encore au besoin $X$
par un voisinage de $y$, on peut supposer que
$f_i\in\Gamma(X,\mathcal J)$ pour $1\leq i\leq p$. En outre, si l'on pose
$\mathcal G=\mathcal F/(\sum_{i=1}^{p}f_i\mathcal F)$, l'hypothèse
$\operatorname{prof}(\mathcal G_y)=0$ entraîne que $\mathcal G_y$ contient
un sous-module isomorphe à
$\mathcal O_{X,y}/\mathfrak m_y$ (0, 16.4.6); le même raisonnement (compte
tenu de (0$_{\rm I}$, 5.3.8 et 5.2.7)) montre qu'on peut supposer qu'il
existe un sous-$\mathcal O_X$-Module $\mathcal G'$ de $\mathcal G$
isomorphe à $\mathcal O_X/\mathcal J$ de sorte que
$\mathcal G'_x=\mathcal O_{Y,x}$ pour tout $x\in Y$. Posons
$\mathcal G''=\mathcal G/\mathcal G'$; remplaçant $X$ par un voisinage
ouvert de $y$, on peut, en vertu de (6.10.2), supposer que $\mathcal G'$ et
$\mathcal G''$ sont normalement plats le long de $Y$. Soit alors $x$ un
point quelconque de $Y$, et posons
$q=\operatorname{prof}(\mathcal O_{Y,x})=\operatorname{prof}(\mathcal G'_x)$;
soit $(g_j)_{1\leq j\leq q}$ une suite maximale
$\mathcal G'_x$-régulière d'éléments de l'idéal maximal
$\mathfrak m_x$ de $\mathcal O_{X,x}$. Chacun des composants homogènes de
$\operatorname{gr}_{\mathcal J_x}^{\bullet}(\mathcal G''_x)$ est un
$\mathcal O_{Y,x}$-module plat de type fini par hypothèse, donc est un
$\mathcal O_{Y,x}$-module libre, et il en est par suite de même de
$\operatorname{gr}_{\mathcal J_x}^{\bullet}(\mathcal G''_x)$; comme la
suite $(g_j)$ est $\mathcal O_{Y,x}$-régulière, elle est par suite
$\operatorname{gr}_{\mathcal J_x}^{\bullet}(\mathcal G''_x)$-régulière,
d'où l'on conclut qu'elle est $\mathcal G''_x$-régulière (0, 15.1.19).
Appliquant le lemme (0, 15.1.18) à la suite exacte
\[
0\longrightarrow\mathcal G'_x\longrightarrow\mathcal G_x
\longrightarrow\mathcal G''_x\longrightarrow0
\]
par récurrence sur $j$, on en conclut une suite exacte
\[
0\longrightarrow\mathcal G'_x/
\left(\sum_{j=1}^{q}g_j\mathcal G'_x\right)
\longrightarrow\mathcal G_x/
\left(\sum_{j=1}^{q}g_j\mathcal G_x\right).
\]
Mais par hypothèse
$\operatorname{prof}(\mathcal G'_x/(\sum_{j=1}^{q}g_j\mathcal G'_x))=0$
(0, 16.4.6); par la caractérisation (0, 16.4.6) des modules de profondeur
nulle, on en conclut que
\[
\operatorname{prof}\left(\mathcal G_x/
\left(\sum_{j=1}^{q}g_j\mathcal G_x\right)\right)=0
=\operatorname{prof}\left(\mathcal F_x/
\left(\sum_{i=1}^{p}f_i\mathcal F_x+
\sum_{j=1}^{q}g_j\mathcal F_x\right)\right);
\]
la suite $(g_j)$ étant $\mathcal G'_x$-régulière et
$\mathcal G''_x$-régulière, est aussi $\mathcal G_x$-régulière; enfin, la
suite $((f_i)_x)$ est $\mathcal F_x$-régulière par hypothèse et est formée
d'éléments de l'idéal maximal $\mathfrak m_x$; on en déduit (0, 16.4.6)
que $\operatorname{prof}(\mathcal F_x)=p+q$, ce qui achève la
démonstration.
\end{proof}

\subsection{\texorpdfstring{Critères pour que les ensembles
$U_{S_n}(\mathcal F)$ ou $U_{C_n}(\mathcal F)$ soient ouverts}{Critères pour que les ensembles U\string_\string{S\string_n\string}(F) ou U\string_\string{C\string_n\string}(F) soient ouverts}}
\label{IV.2.s6.11-fr}

\begin{lemma}[6.11.1]
\label{IV.2.s6.11.1-fr}
Soit $X$ un préschéma localement noethérien, $\mathcal F$ un
$\mathcal O_X$-Module cohérent; alors la fonction
$x\mapsto\operatorname{dim.proj}(\mathcal F_x)$ est semi-continue
supérieurement dans $X$.
\end{lemma}

\begin{proof}
On peut se borner au cas où $X=\operatorname{Spec}(A)$ est le spectre d'un
anneau noethérien et $\mathcal F=\widetilde M$, où $M$ est un $A$-module de
type fini. Supposons que pour un $x\in X$, on ait
$\operatorname{dim.proj}(M_x)=n<+\infty$ (si $n=+\infty$ il n'y a rien à
démontrer); il existe une résolution de $M$
\[
L_{n-1}\longrightarrow L_{n-2}\longrightarrow\cdots
\longrightarrow L_0\longrightarrow M\longrightarrow0
\]

